The Gaussian--Bonferroni diagnostic uses looks \(\lceil20\cdot1.5^j\rceil\) up to \(N\), critical value \(\Phi^{-1}(1-\alpha/(2K))\), sample standard deviation and a finite-population variance correction. It applies the same three labels and returns exact truth at census. Marginal Gaussian pivots are approximate for arbitrary finite vectors; multiplicity correction alone does not make them exact. Certified methods inspect every item, so this is not an inspection-frequency ablation.

This is not Arviv et al.'s Pocock/gsDesign engine. Their source rejects nonpositive aggregate variance; this diagnostic permits zero-width intervals. A failure here cannot be attributed to their implementation.

The Gaussian diagnostic applies multiplicity correction at predeclared looks but retains an approximate sample-variance pivot. On a vector with one 1 and \(N\)-1 zeros, using \(\epsilon\)=.5/\(N\), its decision errors are 292/300 at \(N=1000\) and 299/300 at \(N=5000\), at mean cost fractions .046 and .007. Both certified methods reach census on these runs with no errors; the two-sided 95\% upper binomial limit for zero errors in 300 runs is .01222. This example diagnoses unsupported marginal calibration when rare outcomes have not yet appeared, not optional peeking alone.

The original simulations distinguish full-path miscoverage from final decision error. The former assesses the confidence sequence; the latter assesses the returned label. Neither empirical quantity replaces a uniform guarantee.

\begin{figure}[H]
\centering
\includegraphics[width=.7\textwidth]{figures/gaussian_diagnostic.pdf}
\caption{Rare-event diagnostic: mean cost versus error frequency for two fixed populations (\(N=1000\) and 5000), with 300 permutations per method. Coincident certified-method points overlap. The horizontal line marks \(\alpha=.05\). This panel diagnoses marginal calibration and is separate from the central cost evidence.}\label{fig:gaussian}
\end{figure}


